% Full corrected source proof block; not a Paper II manuscript.
% Checker: GPT-6 Astra (OpenAI, ChatGPT Pro mode), effort Pro.
% Date: 17 November 2026 UTC.
\begin{statement}{E6}{Proved}{Stability of a rank-two torus fibration}
Let $g:Y\to\mathbb P^1$ be as in E5 and suppose its
Kodaira--Spencer tensor has rank two at one regular value. For every
proper deformation $\mathcal Y\to(S,0)$ of $Y$ over a complex
analytic germ, including a nonreduced germ, there is, after shrinking
$S$, a holomorphic map
\[
 G:\mathcal Y\longrightarrow\mathbb P^1\times S
\]
over $S$ whose central restriction is $g$. Every $G_s$ is surjective
with connected fibres and smooth general fibre a complex two-torus.
With $b_2(Y)=0$, every sufficiently small deformation has algebraic
dimension one.
\end{statement}
\begin{proof}
We prove the lifting assertion at every Artin order; this is more
than the tangent-space assertion of E5. The obstruction calculations
below use the Cech construction of holomorphic-map deformations
in \cite[\S4, equations (4.3)--(4.5), and \S6]{Horikawa1974};
we spell out the portions needed here. The injectivity hypothesis
on degree-two tangent cohomology in its Theorem~6.1 is replaced
here by the explicitly proved unobstructedness of the torus source.

First fix a rank-two regular fibre and its embedding $i:F\hookrightarrow
Y$. The normal line is canonically $N=\mathcal O_F\otimes T_c\mathbb
P^1$. The tangent sequence and E5 imply that
\begin{equation}\label{eq:E-surj-embedding}
 H^1(F,T_F)\longrightarrow H^1(F,i^*T_Y)
 \quad\hbox{is surjective}.
\end{equation}
Let $A'\to A$ be a small extension of local Artin complex algebras,
with kernel $I$ annihilated by the maximal ideal. Suppose an embedding
$i_A:F_A\hookrightarrow Y_A$ deforming $i$ and a deformation
$Y_{A'}$ lifting $Y_A$ have been specified. The deformation functor
of a complex torus is smooth: with a real lattice marking its small
complex structures are graphs of arbitrary small complex linear maps
between the two complementary Hodge subspaces; constant Beltrami
tensors have zero bracket. This gives a smooth four-dimensional
miniversal germ for $F$. In particular $F_A$ has a lift $F_{A'}$.

Choose local lifts of $i_A$ on sufficiently small coordinate charts.
They exist because the target deformation is smooth over $A'$.
Their differences on overlaps form a one-cocycle in
$i^*T_Y\otimes I$. Changing the lifted transition functions of
$F_{A'}$ by a cocycle $v_{ij}\in T_F\otimes I$ changes this
obstruction by $di(v_{ij})$. Such changes are precisely the action
of $H^1(F,T_F)\otimes I$ on the set of lifts of $F_A$; their
cocycle law has no quadratic term because $I\mathfrak m_{A'}=0$.
Surjectivity \eqref{eq:E-surj-embedding} therefore makes the map
obstruction zero. A zero-coboundary correction glues the local maps.
The resulting proper map is a closed immersion: the map on local
rings is surjective modulo the nilpotent ideal, hence surjective
by Nakayama, and its underlying central map is the embedding $i$.
Thus the given embedded torus lifts in every ambient small extension.

Now suppose $g_A:Y_A\to\mathbb P^1_A$ has been lifted and that
$Y_{A'}$ is prescribed. Choose local lifts $t_i$ of the target
coordinate of $g_A$; on overlaps their discrepancies define the
holomorphic-map obstruction
\[
 \alpha\in H^1(Y,g^*T_{\mathbb P^1})\otimes I.
\]
Coordinate changes on the target give exactly the transition factors
of $g^*T_{\mathbb P^1}$, and the triple-overlap relation gives its
Cech cocycle law. A coboundary makes the local maps agree, so
$\alpha=0$ is sufficient as well as necessary.

For a rank-two regular value $c$, the divisor
$F_A=g_A^{-1}(c)$ is a smooth deformation of the central torus.
The preceding embedding argument provides a Cartier divisor
$F_{A'}\subset Y_{A'}$ lifting it. This forces the restriction
of $\alpha$ to $F$ to vanish, as can be checked without a comparison
of obstruction theories: choose local equations $t_i-\eta_i$ for
$F_{A'}$, where $\eta_i$ lies in $I\otimes\mathcal O_Y$.
Two such equations differ by a unit. Modulo the central equation
$t=0$, that unit difference disappears, and their overlap equality
reads
\[
 (t_i-t_j)|_F=\eta_i|_F-\eta_j|_F.
\]
This is exactly the assertion that the restricted obstruction
cocycle is a coboundary in $N\otimes I$.

By \eqref{eq:E-Leray-map} and $H^1(\mathbb P^1,T_{\mathbb P^1})=0$,
$\alpha$ is a section of the locally free sheaf
$T_{\mathbb P^1}\otimes R^1g_*\mathcal O_Y$, tensored with $I$.
Its evaluation at every rank-two regular fibre is zero. That set
is open dense; hence the section and $\alpha$ are zero. The map
therefore lifts across every small extension.

We pass from the formal assertion to an analytic map and
specify the relative parameter space, including its nilpotents.
Let $H$ be the germ at the graph of $g$ in the absolute Douady
space of $\mathcal Y\times\mathbb P^1$, and let
$\pi:\mathcal R\to H$ be its universal compact flat family.
Douady's existence and finiteness theorems
\cite[Section~9, Theorems~1--2, pp.~83--85]{Douady1966}
apply to this analytic ambient space; its local rings and those
of the parameter space may have nilpotents.

After shrinking $H$, one has
\[
 \pi_*\mathcal O_{\mathcal R}=\mathcal O_H,
\]
compatibly with base change. To check this without assuming a
reduced parameter space, represent $R\pi_*\mathcal O_{\mathcal R}$
near the central point by a finite complex of finite free analytic
modules. On the central fibre the unit spans
$H^0(Y,\mathcal O_Y)=\mathbb C$. Split off this unit in degree
zero. The differential is injective modulo the maximal ideal on
the complementary degree-zero summand; an invertible minor
cancels that summand with its image in degree one. The resulting
complex begins with $\mathcal O_H$ in degree zero and zero
differential from it. This proves the displayed equality and its
base-change assertion on the full analytic germ.

Embed a representative of $(S,0)$ in a polydisc. The coordinate
functions of the map $\mathcal R\to\mathcal Y\to S$ are now
pullbacks of unique holomorphic functions on $H$. They satisfy
the defining equations of $S$, because their pullbacks do and the
unit map just proved is injective. They therefore define a
holomorphic morphism $H\to S$, and the universal subspace lies
scheme-theoretically in
\[
 (\mathcal Y\times_S H)\times\mathbb P^1.
\]
The projection of this subspace to $\mathcal Y\times_S H$ is a
proper morphism between two proper flat families over $H$ and
is an isomorphism at the central point. By the isomorphism-locus
criterion \cite[Section~10.1, Proposition~1, p.~86]{Douady1966},
it is an isomorphism after restricting to an open neighbourhood
of that point. The proof of that criterion uses the finite
direct image, flatness and Nakayama's lemma, and applies to
nonreduced $H$. On this open locus the universal subspace is
the graph of an actual map to $\mathbb P^1$.

The universal property of the absolute Douady space consequently
identifies this open locus with the germ $\mathcal H\to S$ of
relative maps from $\mathcal Y/S$ to
$\mathbb P^1\times S/S$. This is the relative version of
Douady's graph construction
\cite[Section~10.2, Theorem~1, p.~87]{Douady1966}; the argument
above supplies the variation of the compact source explicitly.
The just-proved lifting property for every
small Artin extension says that this morphism is formally smooth
at $[g]$. The local analytic smoothness criterion identifies a
formally smooth morphism of finite-dimensional analytic germs with
a projection $S\times(\mathbb C^r,0)\to S$ after shrinking. One
can see the criterion on presentations of their convergent local
algebras: lifting across all square-zero quotients makes the
relative conormal obstruction zero, and the analytic implicit
function theorem eliminates the corresponding equations. The
zero section of this product supplies the required holomorphic
family of graphs. This argument applies to nonreduced $S$.

At one regular value, the family $G$ remains a proper submersion
near the central compact fibre. Its nearby fibres are connected
deformations of a torus. Each map $G_s$ is nonconstant, and hence
surjective. Its Stein factorization has finite degree one because
it has a connected regular fibre in this neighbourhood. Thus
every fibre is connected. Over the connected smooth locus of
$G_s$, Ehresmann's theorem identifies the differentiable fibres
with $T^4$. Every compact complex surface with this differentiable
type is a torus: its even first Betti number makes it K\"ahler
by \cite[Theorem 11]{Buchdahl1999};
its Albanese has dimension two; the exterior algebra on its four
integral degree-one classes shows the Albanese map has degree one;
and a nontrivial exceptional curve of this degree-one map would
include a $(-1)$-curve, impossible for the even intersection form
of $T^4$. Hence the general fibre is a complex two-torus.
Finally $b_2$ is constant in the differentiable deformation family,
and E1 proves $a(Y_s)=1$.
\end{proof}
